\chapter{Quasi-Free States, Wick's Theorem, and the Variational Grassmannian}
\label{ch:m1_1b}

% EDITORIAL QUESTION: Why is Wick factorisation a property of the state rather than an operator identity?
\section{Motivating Example: The Failure of Operator Decoupling on Two Spins}
\label{sec:m1_1b_motivating_example}
In conventional treatments of condensed matter mean-field theory, Hartree--Fock--Bogoliubov decoupling is frequently presented as an intuitive operator replacement:
\begin{equation}
\label{eq:m1_1b_naive_decoupling}
A B C D \stackrel{?}{\approx} \langle A B \rangle C D + A B \langle C D \rangle - \langle A C \rangle B D - \dots
\end{equation}
Such an expression cannot hold as an algebraic identity between operators.

% CLAIM: CLM-M1-QF-001 (provenance: doc/rnd/phase_1/m1_5_t3_quasifree_states_wick_variational.md, test: test/unit/test_quasifree_states.py::test_tst_qf_02_wick_fails_for_generic_non_gaussian_state)
\begin{proposition}[Failure of Operator-Level Wick Decoupling on States]
\label{prop:m1_1b_wick_failure}
Consider two spin-$1/2$ moments on a single bond $\langle ij \rangle$ governed by the ferromagnetic Kitaev interaction $H = -K \sigma_i^z \sigma_j^z$ with $K > 0$. Using the exact Majorana representation $\sigma_i^z = -\ii c_i^0 c_i^z$ and $\sigma_j^z = -\ii c_j^0 c_j^z$, the exchange Hamiltonian is quartic:
\begin{equation}
\label{eq:m1_1b_two_spin_quartic}
H = K (c_i^0 c_i^z)(c_j^0 c_j^z).
\end{equation}
For an entangled spin singlet state $|\Psi_{\mathrm{singlet}}\rangle = \frac{1}{\sqrt{2}}(|\!\uparrow\downarrow\rangle - |\!\downarrow\uparrow\rangle)$, the truncated fourth cumulant:
\begin{equation}
\label{eq:m1_1b_fourth_cumulant}
\kappa_4(c_i^0, c_i^z, c_j^0, c_j^z) \equiv \langle c_i^0 c_i^z c_j^0 c_j^z \rangle - \left[ \langle c_i^0 c_i^z \rangle \langle c_j^0 c_j^z \rangle - \langle c_i^0 c_j^0 \rangle \langle c_i^z c_j^z \rangle + \langle c_i^0 c_j^z \rangle \langle c_i^z c_j^0 \rangle \right],
\end{equation}
does not vanish ($\kappa_4 \neq 0$).
\end{proposition}
\begin{proof}
Direct evaluation on the 4-dimensional Fock space shows that the exact singlet state has $\langle \sigma_i^z \sigma_j^z \rangle = -1$. In contrast, the product of two-point expectations evaluates to zero or positive values depending on gauge alignment. Thus $\kappa_4 \ne 0$ identically.
\end{proof}

This elementary counterexample establishes the epistemological foundation of this chapter: \textbf{Wick's theorem is not an identity on the operator algebra; it is a structural property characterizing a restricted family of states---the quasi-free (or Gaussian) states.} Mean-field theory is the restriction of the quantum variational principle to this state manifold.

% EDITORIAL QUESTION: How is a finite spin system formulated rigorously in algebraic quantum field theory (AQFT)?
\section{Finite Quasi-Local Spin Nets and the State Space}
\label{sec:m1_1b_aqft_spin_nets}
% SOURCE [VERIFIED]: Bratteli & Robinson (1997), Operator Algebras and Quantum Statistical Mechanics II, Sec. 5.2 — Quasi-local spin systems
% CLAIM: CLM-M1-QF-002 (provenance: doc/rnd/phase_1/m1_5_t3_quasifree_states_wick_variational.md)
Let $\Lambda = (\mathcal{V}, \mathcal{E})$ be a finite bipartite graph consisting of $N = |\mathcal{V}|$ lattice sites.

\begin{definition}[Finite Quasi-Local Spin Net]
\label{def:m1_1b_spin_net}
To each subset of sites $X \subseteq \mathcal{V}$, associate the finite-dimensional $C^*$-algebra $\mathcal{A}_X \equiv \bigotimes_{i \in X} \mathcal{B}(\mathcal{H}_i) \cong \bigotimes_{i \in X} M_2(\mathbb{C})$. The family $\{\mathcal{A}_X\}_{X \subseteq \mathcal{V}}$ constitutes a quasi-local net satisfying:
\begin{enumerate}
\item \textbf{Isotony:} $X \subseteq Y \implies \mathcal{A}_X \subseteq \mathcal{A}_Y$ via the canonical unital embedding $A \mapsto A \otimes \Id_{Y \setminus X}$.
\item \textbf{Microcausality (Strict Locality):} If $X \cap Y = \emptyset$, then $[A, B] = 0$ for all $A \in \mathcal{A}_X$ and $B \in \mathcal{A}_Y$.
\item \textbf{Quasi-Local Observable Algebra:} The global observable algebra is $\ALambda \equiv \mathcal{A}_{\mathcal{V}} \cong M_{2^N}(\mathbb{C})$.
\end{enumerate}
\end{definition}

A \textbf{state} $\omega$ on $\ALambda$ is a normalized positive linear functional $\omega: \ALambda \to \mathbb{C}$ ($\omega(\Id) = 1$ and $\omega(A^\dagger A) \ge 0$). By the finite-dimensional Gelfand--Naimark--Segal (GNS) theorem, every state is uniquely represented by a density matrix $\rho \in \ALambda$ such that:
\begin{equation}
\label{eq:m1_1b_density_matrix}
\rho \ge 0, \qquad \Tr(\rho) = 1, \qquad \omega(A) = \Tr(\rho A).
\end{equation}
The state space $\mathcal{S}(\ALambda)$ is convex and compact. A local Hamiltonian is specified by an interaction potential $\Phi: \mathcal{P}(\mathcal{V}) \to \ALambda$ with $\Phi(X) = \Phi(X)^\dagger \in \mathcal{A}_X$, yielding:
\begin{equation}
\label{eq:m1_1b_hamiltonian_general}
H = \sum_{X \subseteq \mathcal{V}} \Phi(X).
\end{equation}
The exact quantum ground-state energy is the convex minimum over all states:
\begin{equation}
\label{eq:m1_1b_exact_variational_energy}
E_0 = \min_{\omega \in \mathcal{S}(\ALambda)} \omega(H) = \min_{\substack{\rho \ge 0 \\ \Tr\rho = 1}} \Tr(\rho H).
\end{equation}

% EDITORIAL QUESTION: How does the spin net embed into the self-dual CAR algebra, and what is the role of parity superselection?
\section{Majorana Representation Theory and Graded Locality}
\label{sec:m1_1b_representation_theory}
% CLAIM: CLM-M1-QF-003 (provenance: doc/rnd/phase_1/m1_5_t3_quasifree_states_wick_variational.md, test: test/unit/test_quasifree_states.py::test_tst_qf_08_graded_locality_vs_spin_commutativity)
Let $V \cong \mathbb{R}^{2n}$ with $2n = 8N$ be the real Euclidean space spanned by the Majorana generators.

\begin{definition}[Self-Dual CAR Algebra and $\mathbb{Z}_2$-Grading]
\label{def:m1_1b_car_grading}
The self-dual Canonical Anticommutation Relations algebra $\mathcal{A}_{\mathrm{CAR}}(V) \cong \Cl(V)$ is the unital $*$-algebra generated by Hermitian elements $\{c_a : a = 1, \dots, 2n\}$ satisfying $\{c_a, c_b\} = 2\delta_{ab}\Id$. The parity automorphism $\Theta(c_a) = -c_a$ induces the canonical $\mathbb{Z}_2$-grading:
\begin{equation}
\label{eq:m1_1b_z2_grading}
\mathcal{A}_{\mathrm{CAR}} = \mathcal{A}_{\mathrm{CAR}}^+ \oplus \mathcal{A}_{\mathrm{CAR}}^-,
\end{equation}
where even elements commute at spacelike separation, while odd elements anticommute (\textbf{graded locality}):
\begin{equation}
\label{eq:m1_1b_graded_locality}
A \in \mathcal{A}_X^-, \quad B \in \mathcal{A}_Y^- \implies \{A, B\} = 0 \quad (X \cap Y = \emptyset).
\end{equation}
\end{definition}

\begin{theorem}[Parity Superselection and Spin Embedding]
\label{thm:m1_1b_parity_superselection}
The local spin algebra $\mathcal{A}_{\mathrm{spin}, i}$ embeds into the parity-even gauge-invariant subalgebra of $\mathcal{A}_{\mathrm{CAR}, i}$:
\begin{equation}
\label{eq:m1_1b_spin_embedding}
\mathcal{A}_{\mathrm{spin}, i} \cong P_{G, i} \mathcal{A}_{\mathrm{CAR}, i}^+ P_{G, i}, \qquad P_{G, i} = \frac{\Id + \mathcal{D}_i}{2}.
\end{equation}
Because physical spin operators are even quadratic monomials ($\sigma_i^\alpha = -\ii c_i^0 c_i^\alpha \in \mathcal{A}_{\mathrm{CAR}, i}^+$), they commute identically on distinct lattice sites:
\begin{equation}
\label{eq:m1_1b_spin_commutativity}
[\sigma_i^\alpha, \sigma_j^\beta] = 0 \quad \forall i \ne j,
\end{equation}
recovering ordinary bosonic microcausality despite the underlying fermionic CAR grading.
\end{theorem}

% EDITORIAL QUESTION: How are quasi-free states defined on the CAR algebra, and what makes Wick's theorem exact?
\section{Quasi-Free States and the Pfaffian Theorem}
\label{sec:m1_1b_quasifree_states}
% SOURCE [VERIFIED]: Araki (1970), Publ. RIMS Kyoto Univ. 6, 385-442 — On quasifree states of the CAR algebra
% CLAIM: CLM-M1-QF-004 (provenance: doc/rnd/phase_1/m1_5_t3_quasifree_states_wick_variational.md, test: test/unit/test_quasifree_states.py::test_tst_qf_01_pfaffian_wick_for_gaussian_states)

\begin{definition}[Quasi-Free State]
\label{def:m1_1b_quasifree_state}
A state $\omega$ on $\mathcal{A}_{\mathrm{CAR}}(V)$ is called \textbf{quasi-free (or Gaussian)} if all odd correlation functions vanish, $\omega(c_{a_1} \cdots c_{a_{2k+1}}) = 0$, and all even correlation functions evaluate to the Pfaffian of the two-point function:
\begin{equation}
\label{eq:m1_1b_pfaffian_wick}
\omega(c_{a_1} c_{a_2} \cdots c_{a_{2k}}) = \ii^k \operatorname{Pf}\left[ \Omega_{a_p a_q} \right]_{1 \le p < q \le 2k},
\end{equation}
where $\Omega \in \mathfrak{so}(2n)$ is the real skew-symmetric covariance matrix:
\begin{equation}
\label{eq:m1_1b_covariance_def}
\omega(c_a c_b) = \delta_{ab}\Id + \ii \Omega_{ab}, \qquad \Omega^T = -\Omega \in \mathbb{R}^{2n \times 2n}.
\end{equation}
\end{definition}

\begin{theorem}[Quantum Admissibility and Purity]
\label{thm:m1_1b_admissibility}
A skew-symmetric matrix $\Omega \in \mathfrak{so}(2n)$ specifies a physical quantum state on $\mathcal{A}_{\mathrm{CAR}}$ if and only if its operator norm satisfies:
\begin{equation}
\label{eq:m1_1b_admissibility_bound}
\|\Omega\|_2 \le 1 \iff -\Id \le \Omega^2 \le 0.
\end{equation}
The state is pure if and only if $\Omega$ is an orthogonal complex structure:
\begin{equation}
\label{eq:m1_1b_purity_condition}
\Omega^2 = -\Id \iff \operatorname{spec}(\Omega) = \{\pm \ii\}.
\end{equation}
For any pure quasi-free state, $\operatorname{Pf}(\Omega) = \pm 1$ represents the global fermion parity.
\end{theorem}
\begin{proof}
Let $\omega$ be a state on $\mathcal{A}_{\mathrm{CAR}}$. Positivity requires that for any complex vector $v \in \mathbb{C}^{2n}$, the element $A = \sum_a v_a c_a$ satisfies $\omega(A^\dagger A) \ge 0$. Direct evaluation gives $\omega(A^\dagger A) = v^\dagger (\Id + \ii\Omega) v \ge 0$, which is equivalent to the matrix operator inequality $\Id + \ii\Omega \succeq 0$. Since $\ii\Omega$ is Hermitian with eigenvalues $\pm \lambda_k$ ($\lambda_k \in \mathbb{R}$), $1 \pm \lambda_k \ge 0$ requires $|\lambda_k| \le 1$, proving Eq.~\eqref{eq:m1_1b_admissibility_bound}. Purity corresponds to extremal states in the convex set, which occurs if and only if all eigenvalues saturate the boundary $|\lambda_k| = 1$, i.e., $\Omega^2 = -\Id$.
\end{proof}

% EDITORIAL QUESTION: What is the differential geometry of the pure quasi-free state manifold?
\section{Differential Geometry of the Real Orthogonal Grassmannian}
\label{sec:m1_1b_grassmannian_geometry}
% SOURCE [VERIFIED]: Bach, Lieb & Solovej (1994), J. Stat. Phys. 76, 3-89 — Generalized Hartree-Fock theory
% CLAIM: CLM-M1-QF-005 (provenance: doc/rnd/phase_1/m1_5_t3_quasifree_states_wick_variational.md, test: test/unit/test_quasifree_states.py::test_tst_qf_04_admissibility_purity_and_fermion_parity)

\begin{definition}[The Real Orthogonal Grassmannian $\Gr(n, 2n)$]
\label{def:m1_1b_grassmannian}
The manifold of pure quasi-free states on $2n$ Majorana modes is isomorphic to the compact homogeneous space:
\begin{equation}
\label{eq:m1_1b_grassmannian_manifold}
\mathcal{M}_{\mathrm{pure}} \cong \frac{\Oo(2n)}{\U(n)} \cong \Gr(n, 2n).
\end{equation}
Its real dimension is $\dim_{\mathbb{R}} \Gr(n, 2n) = n(n-1)$. It consists of two disconnected components distinguished by the parity invariant $\operatorname{Pf}(\Omega) \in \{+1, -1\}$.
\end{definition}

Each pure covariance $\Omega \in \Gr(n, 2n)$ uniquely defines an orthogonal projector $P$ onto a maximal isotropic subspace of $V \otimes \mathbb{C}$:
\begin{equation}
\label{eq:m1_1b_projector_covariance}
P = \frac{\Id + \ii \Omega}{2}, \qquad P^2 = P = P^\dagger, \quad \Tr(P) = n.
\end{equation}
The tangent space to the Grassmannian at $\Omega$ is given by:
\begin{equation}
\label{eq:m1_1b_tangent_space}
T_\Omega \Gr(n, 2n) = \{ X \in \mathfrak{so}(2n) : X \Omega + \Omega X = 0 \}.
\end{equation}

\begin{proposition}[Non-Convexity of the Quasi-Free Manifold]
\label{prop:m1_1b_non_convexity}
While the set of admissible covariance matrices $\{\Omega \in \mathfrak{so}(2n) : \|\Omega\|_2 \le 1\}$ is convex, the corresponding set of quasi-free states $\mathcal{M}_{\mathrm{QF}} \subset \mathcal{S}(\ALambda)$ is \textbf{not convex}. A statistical mixture of two Gaussian states $\rho_{\mathrm{mix}} = \frac{1}{2}\rho_1 + \frac{1}{2}\rho_2$ possesses a non-vanishing fourth cumulant $\kappa_4 \ne 0$.
\end{proposition}
\begin{proof}
Let $\rho_1, \rho_2 \in \mathcal{M}_{\mathrm{QF}}$ be distinct Gaussian density matrices with covariances $\Omega_1 \neq \Omega_2$. The mixture $\rho_{\mathrm{mix}} = \frac{1}{2}(\rho_1 + \rho_2)$ has two-point covariance $\Omega_{\mathrm{mix}} = \frac{1}{2}(\Omega_1 + \Omega_2)$. Computing the four-point expectation $\Tr(\rho_{\mathrm{mix}} c_1 c_2 c_3 c_4) = \frac{1}{2}\operatorname{Pf}(\Omega_1) + \frac{1}{2}\operatorname{Pf}(\Omega_2)$, whereas a Gaussian state with covariance $\Omega_{\mathrm{mix}}$ would have $\operatorname{Pf}(\Omega_{\mathrm{mix}}) = \operatorname{Pf}\left(\frac{\Omega_1 + \Omega_2}{2}\right)$. Because the Pfaffian is a non-linear polynomial of degree 2 in $4 \times 4$ matrices, $\frac{1}{2}\operatorname{Pf}(\Omega_1) + \frac{1}{2}\operatorname{Pf}(\Omega_2) \neq \operatorname{Pf}\left(\frac{\Omega_1 + \Omega_2}{2}\right)$ in general, yielding non-zero fourth cumulant $\kappa_4 \ne 0$.
\end{proof}

\begin{figure}[htbp]
\centering
\begin{tikzpicture}[scale=1.05]
    % Canvas background
    \draw[thick, fill=labbookbg!80!white, draw=labbookrule] (-6.8, -3.0) rectangle (6.8, 3.2);
    \node[anchor=north west, font=\normalsize\bfseries\sffamily, color=labbookslate] at (-6.5, 2.95) {Total Quantum State Space $\mathcal{S}(\ALambda)$ (Convex)};

    % Physical Subspace (Gutzwiller Projection hyperplane)
    \draw[very thick, labbookteal, dashed] (-5.8, 1.6) -- (3.2, 1.6);
    \node[anchor=west, font=\scriptsize\bfseries, color=labbookteal, align=left] at (3.35, 1.6) {Physical Subspace\\$\mathcal{F}_{\mathrm{phys}} = \Ran P_G$ ($\mathcal{D}_i = +\Id$)};

    % Non-Convex Quasi-Free Submanifold M_QF (Shaded Region)
    \draw[very thick, fill=labbookwarm!30!white, draw=labbookcrimson]
        (-4.6, -2.2) [rounded corners=20pt] -- (-2.4, -0.3) [rounded corners=18pt] -- (0.0, -0.9) [rounded corners=18pt] -- (2.0, -0.3) [rounded corners=20pt] -- (4.2, -2.2) [rounded corners=20pt] -- (3.6, -2.4) [rounded corners=20pt] -- (-4.0, -2.4) -- cycle;
    \node[anchor=center, font=\small\bfseries, color=labbookcrimson] at (-0.5, -1.8) {Quasi-Free States $\mathcal{M}_{\mathrm{QF}}$ (Non-Convex)};

    % Pure States: Orthogonal Grassmannian Gr(n, 2n) Boundary
    \draw[very thick, labbookaccent]
        (-4.6, -2.2) to[out=45, in=185] (-2.4, -0.3) to[out=5, in=175] (0.0, -0.9) to[out=-5, in=185] (2.0, -0.3) to[out=5, in=135] (4.2, -2.2);
    \node[anchor=south west, font=\footnotesize\bfseries, color=labbookaccent, align=left] at (-6.2, -0.7) {Grassmannian $\Gr(n, 2n)$\\($\Omega^2 = -\Id$)};

    % State points on Physical Hyperplane
    \filldraw[labbookteal] (-2.5, 1.6) circle (3.5pt) node[anchor=south, font=\small\bfseries, color=labbookteal, yshift=3pt] {$\rho_{\mathrm{exact}}$};
    \filldraw[labbookcrimson] (2.0, 1.6) circle (3.5pt) node[anchor=south, font=\small\bfseries, color=labbookcrimson, yshift=3pt] {$P_G \Omega_{\mathrm{MF}} P_G$};

    % State point on Grassmannian
    \filldraw[labbookaccent] (2.0, -0.3) circle (3.5pt) node[anchor=west, font=\small\bfseries, color=labbookaccent, xshift=6pt] {$\Omega_{\mathrm{MF}} \in \Gr$};

    % Projection Gap Vector (Vertical Arrow)
    \draw[->, very thick, shorten >=5pt, shorten <=5pt, color=labbookslate] (2.0, -0.3) -- (2.0, 1.6);
    \node[anchor=west, font=\scriptsize\bfseries, color=labbookslate, align=left] at (2.15, 0.75) {Gutzwiller Projection\\Gap (BL-01)};

    % Non-Gaussianity Relative Entropy Distance (Dashed Diagonal Arrow)
    \draw[<->, very thick, dashed, shorten >=6pt, shorten <=6pt, color=labbookcrimson] (-2.5, 1.6) -- (2.0, -0.3)
        node[pos=0.42, below=5pt, sloped, font=\scriptsize\bfseries, color=labbookcrimson] {Variational Gap $D_{\mathrm{NG}}(\rho)$ (BL-07)};
\end{tikzpicture}
\caption{Geometry of the variational quantum state space. The full density matrix space $\mathcal{S}(\ALambda)$ is convex. Quasi-free states $\mathcal{M}_{\mathrm{QF}}$ form a non-convex embedded submanifold whose pure boundary is the real orthogonal Grassmannian $\Gr(n, 2n)$. The mean-field state $\Omega_{\mathrm{MF}}$ minimizes energy on the Grassmannian, while Gutzwiller projection maps onto the physical gauge-singlet hyperplane.}
\label{fig:m1_1b_grassmannian_schematic}
\end{figure}


% EDITORIAL QUESTION: How does mean-field decoupling emerge as the Euler-Lagrange stationarity condition on the Grassmannian?
\section{The Variational Principle on the Grassmannian and Exact Mean-Field Decoupling}
\label{sec:m1_1b_variational_principle}
% CLAIM: CLM-M1-QF-006 (provenance: doc/rnd/phase_1/m1_5_t3_quasifree_states_wick_variational.md, test: test/unit/test_quasifree_states.py::test_tst_qf_06_affine_tangent_property, test_tst_qf_07_grassmannian_stationarity_commutator)

\begin{definition}[General Quartic Majorana Hamiltonian]
\label{def:m1_1b_general_hamiltonian}
Let $A \in \mathfrak{so}(2n)$ be a real skew-symmetric matrix describing one-body hoppings, and let $V_{abcd}$ be a real totally antisymmetric 4-tensor describing two-body interactions. The general microscopic Hamiltonian is:
\begin{equation}
\label{eq:m1_1b_general_quartic_H}
H = \frac{\ii}{4}\sum_{a,b=1}^{2n} A_{ab} c_a c_b + \frac{1}{16}\sum_{a,b,c,d=1}^{2n} V_{abcd} c_a c_b c_c c_d.
\end{equation}
\end{definition}

\begin{definition}[Hartree-Fock-Bogoliubov Self-Energy and Single-Particle MF Matrix]
\label{def:m1_1b_hfb_self_energy}
For any covariance matrix $\Omega \in \mathfrak{so}(2n)$, the Hartree--Fock--Bogoliubov (HFB) self-energy matrix $\Sigma(\Omega) \in \mathfrak{so}(2n)$ is defined by the contraction:
\begin{equation}
\label{eq:m1_1b_sigma_self_energy}
\Sigma_{ab}(\Omega) \equiv -\frac{1}{2} \sum_{c, d=1}^{2n} V_{abcd} \Omega_{cd}.
\end{equation}
The single-particle mean-field Hamiltonian matrix $h(\Omega) \in \mathfrak{so}(2n)$ is the sum of the bare one-body hopping matrix $A$ and the self-energy $\Sigma(\Omega)$:
\begin{equation}
\label{eq:m1_1b_h_omega_def}
h(\Omega) \equiv A + \Sigma(\Omega).
\end{equation}
\end{definition}

\begin{theorem}[Energy Functional and Gradient on the Grassmannian]
\label{thm:m1_1b_energy_gradient}
Restricting the state expectation $\omega_\Omega(H)$ to the quasi-free manifold via Wick's theorem yields the polynomial energy functional:
\begin{equation}
\label{eq:m1_1b_quasifree_energy}
\mathcal{E}(\Omega) = E_0^{\mathrm{core}} - \frac{1}{4} \Tr(A \Omega) - \frac{1}{8} \sum_{a,b,c,d} V_{abcd} (\Omega_{ab}\Omega_{cd} - \Omega_{ac}\Omega_{bd} + \Omega_{ad}\Omega_{bc}).
\end{equation}
The single-particle mean-field Hamiltonian matrix $h(\Omega) \in \mathfrak{so}(2n)$ is the functional gradient of $\mathcal{E}$ with respect to $\Omega$:
\begin{equation}
\label{eq:m1_1b_hmf_gradient}
h(\Omega) = 4 \frac{\partial \mathcal{E}}{\partial \Omega} = A + \Sigma(\Omega).
\end{equation}
A state $\Omega \in \Gr(n, 2n)$ is stationary on the Grassmannian if and only if the matrix commutator vanishes:
\begin{equation}
\label{eq:m1_1b_euler_lagrange_stationarity}
[h(\Omega), \Omega] = 0 \iff \Omega = -\operatorname{sgn}(\ii h(\Omega)).
\end{equation}
\end{theorem}
\begin{proof}
By Wick's theorem (Definition~\ref{def:m1_1b_quasifree_state}), $\omega_\Omega(c_a c_b) = \delta_{ab} + \ii \Omega_{ab}$, so $\omega_\Omega(\frac{\ii}{4} c^T A c) = -\frac{1}{4}\Tr(A\Omega)$. For the quartic term, $\omega_\Omega(c_a c_b c_c c_d) = -(\Omega_{ab}\Omega_{cd} - \Omega_{ac}\Omega_{bd} + \Omega_{ad}\Omega_{bc})$. Taking the derivative with respect to the skew-symmetric matrix $\Omega_{ab}$ using totally antisymmetric $V_{abcd}$, each contraction yields $4 \frac{\partial \mathcal{E}}{\partial \Omega_{ab}} = A_{ab} + \Sigma_{ab}(\Omega) = h_{ab}(\Omega)$.

To establish stationarity on $\Gr(n, 2n)$, let $\Omega(t) = e^{t X} \Omega e^{-t X}$ be a smooth curve on the Grassmannian generated by tangent vector $X \in T_\Omega \Gr$. Differentiating along the curve:
\begin{equation}
\left.\frac{\dd}{\dd t}\mathcal{E}(\Omega(t))\right|_{t=0} = \frac{1}{4} \Tr\left( h(\Omega)^T [X, \Omega] \right) = \frac{1}{4} \Tr\left( [h(\Omega), \Omega] X \right).
\end{equation}
Stationarity requires $\Tr([h(\Omega), \Omega] X) = 0$ for all tangent vectors $X$. Since $T_\Omega \Gr$ consists of all skew matrices anticommuting with $\Omega$, $[h(\Omega), \Omega]$ lies in $T_\Omega \Gr$, which forces $[h(\Omega), \Omega] = 0$. Since $h(\Omega)$ is skew-symmetric, its non-zero eigenvalues are imaginary conjugate pairs $\pm \ii \epsilon_k$ ($\epsilon_k > 0$). To minimize energy, $\Omega$ must assign eigenvalue $-\ii$ to positive energy modes and $+\ii$ to negative energy modes, yielding $\Omega = -\operatorname{sgn}(\ii h(\Omega))$.
\end{proof}

\begin{theorem}[Affine Tangent Property and Scalar Double-Counting Subtraction]
\label{thm:m1_1b_affine_tangent}
Let $H_{\mathrm{MF}}[\Omega] \equiv \frac{\ii}{4} c^T h(\Omega) c$ be the quadratic mean-field operator. Then:
\begin{equation}
\label{eq:m1_1b_double_counting_scalar}
\omega_\Omega(H_{\mathrm{MF}}[\Omega]) = 2 E_{\mathrm{int}}(\Omega) + E_{\mathrm{quad}}(\Omega).
\end{equation}
The exact variational energy requires subtracting the scalar double-counting correction $C[\Omega] \equiv E_{\mathrm{int}}(\Omega)$:
\begin{equation}
\label{eq:m1_1b_energy_subtraction}
\mathcal{E}(\Omega) = \omega_\Omega(H_{\mathrm{MF}}[\Omega]) - C[\Omega].
\end{equation}
Moreover, $H_{\mathrm{MF}}[\Omega]$ defines an affine tangent plane to the energy surface:
\begin{equation}
\label{eq:m1_1b_affine_tangent_bound}
\mathcal{E}(\Omega') - \left[ \omega_{\Omega'}(H_{\mathrm{MF}}[\Omega]) - C[\Omega] \right] = \mathcal{O}(\|\Omega' - \Omega\|^2).
\end{equation}
\end{theorem}
\begin{proof}
The interaction energy $E_{\mathrm{int}}(\Omega) = -\frac{1}{8}\sum V_{abcd}\Omega_{ab}\Omega_{cd}$ is a quadratic homogeneous polynomial in $\Omega$. By Euler's homogeneous function theorem, $\sum_{ab}\Omega_{ab}\frac{\partial E_{\mathrm{int}}}{\partial \Omega_{ab}} = 2 E_{\mathrm{int}}(\Omega)$. Since $h(\Omega) = A + 4 \frac{\partial E_{\mathrm{int}}}{\partial \Omega}$, contracting with $\Omega$ yields:
\begin{equation}
\omega_\Omega(H_{\mathrm{MF}}[\Omega]) = -\frac{1}{4}\Tr(h(\Omega)\Omega) = -\frac{1}{4}\Tr(A\Omega) - \Tr\left(\frac{\partial E_{\mathrm{int}}}{\partial \Omega}\Omega\right) = E_{\mathrm{quad}}(\Omega) + 2 E_{\mathrm{int}}(\Omega).
\end{equation}
Subtracting $C[\Omega] = E_{\mathrm{int}}(\Omega)$ exactly recovers $\mathcal{E}(\Omega) = E_{\mathrm{quad}}(\Omega) + E_{\mathrm{int}}(\Omega)$, proving Eq.~\eqref{eq:m1_1b_energy_subtraction}. Expanding $\mathcal{E}(\Omega')$ to first order around $\Omega$ gives $\mathcal{E}(\Omega') = \mathcal{E}(\Omega) + \frac{1}{4}\Tr(h(\Omega)(\Omega' - \Omega)) + \mathcal{O}(\|\Omega'-\Omega\|^2)$, which identically matches the linearized mean-field expectation, establishing Eq.~\eqref{eq:m1_1b_affine_tangent_bound}.
\end{proof}

% EDITORIAL QUESTION: How are local gauge constraints treated in the variational framework, and what is the Gutzwiller gap?
\section{Gauge Constraints, KKT Stationarity, and the Projection Gap}
\label{sec:m1_1b_gauge_kkt}
On the physical spin space, the gauge condition $\mathcal{D}_i = \Id$ must hold strictly on every site $i \in \mathcal{V}$. Within a quasi-free state, $\omega_\Omega(\mathcal{D}_i) = \operatorname{Pf}(\Omega_i)$ is a local polynomial constraint. Because a pure Gaussian state is completely fixed by its two-point correlation matrix $\Omega$, demanding $\operatorname{Pf}(\Omega_i) = 1$ on every site simultaneously while allowing non-trivial inter-site entanglement forces the state outside the quasi-free manifold.

Instead, mean-field theory enforces the local gauge conditions \textbf{on average} via Lagrange multipliers $\boldsymbol{\lambda} = \{\lambda_i^\gamma\}$:
\begin{equation}
\label{eq:m1_1b_kkt_lagrangian}
\mathcal{L}(\Omega, \boldsymbol{\lambda}) = \mathcal{E}(\Omega) - \sum_{i \in \mathcal{V}} \sum_{\gamma=1}^3 \lambda_i^\gamma \omega_\Omega(\tau_i^\gamma).
\end{equation}
The Karush--Kuhn--Tucker (KKT) stationarity conditions read:
\begin{equation}
\label{eq:m1_1b_kkt_stationarity}
\nabla_\Omega \mathcal{L} \in (T_\Omega \Gr)^\perp, \qquad g_i^\gamma(\Omega) \equiv \omega_\Omega(\tau_i^\gamma) = 0.
\end{equation}

\begin{definition}[Expectation Value of the Gauge Generator]
\label{def:m1_1b_gauge_expectation}
From Definition~\ref{def:m1_1_spin_gauge_pauli}, the gauge generator on site $i$ is $\tau_i^\gamma = \frac{\ii}{4} c_i^T \mathrm{G}^\gamma c_i$. In any quasi-free state $\omega_\Omega$, its expectation value evaluates to the trace:
\begin{equation}
\label{eq:m1_1b_gauge_expectation_trace}
\omega_\Omega(\tau_i^\gamma) = -\frac{1}{4} \Tr(\mathrm{G}^\gamma \Omega_{ii}) = -\frac{1}{2} (\Omega_{i}^{0\gamma} + \Omega_{i}^{\gamma^+ \gamma^-}).
\end{equation}
Vanishing expectation $\omega_\Omega(\tau_i^\gamma) = 0$ requires that the on-site matter-gauge covariance $\Omega_{i}^{0\gamma}$ exactly cancel the gauge-gauge cross-link covariance $\Omega_{i}^{\gamma^+ \gamma^-}$.
\end{definition}

The discrepancy between the unprojected expectation $\langle \mathcal{O} \rangle_{\mathrm{G}} = \omega_\Omega(\mathcal{O})$ and the physical Gutzwiller-projected expectation $\langle \mathcal{O} \rangle_{\mathrm{phys}} = \frac{\Tr(\rho P_G \mathcal{O} P_G)}{\Tr(\rho P_G)}$ is the \textbf{Gutzwiller Projection Gap} (cataloged as exploratory backlog item \textbf{BL-01}).

% EDITORIAL QUESTION: How does the quasi-free upper bound interface with semidefinite programming (SDP) lower bounds?
\section{Semidefinite Relaxations, N-Representability, and 2-RDM Lower Bounds}
\label{sec:m1_1b_sdp_bounds}
% CLAIM: CLM-M1-QF-008 (provenance: doc/rnd/phase_1/m1_5_t3_quasifree_states_wick_variational.md)
The Grassmannian variational principle yields a strict \textbf{upper bound} to the true ground energy:
\begin{equation}
\label{eq:m1_1b_upper_bound}
E_0 \le \min_{\Omega \in \Gr(n, 2n)} \mathcal{E}(\Omega) = E_{\mathrm{MF}}.
\end{equation}
Conversely, the fermionic $N$-representability problem yields rigorous \textbf{lower bounds} via semidefinite programming (SDP). Let $\rho^{(2)}$ denote the two-body reduced density matrix with matrix elements $\rho^{(2)}_{ab,cd} = \Tr(\rho c_a c_b c_c c_d)$. The fundamental positivity conditions:
\begin{enumerate}
    \item \textbf{D-Condition (Two-Particle Positivity):} $\langle \psi | c_a c_b (c_c c_d)^\dagger | \psi \rangle \ge 0 \implies \rho^{(2)} \succeq 0$.
    \item \textbf{Q-Condition (Two-Hole Positivity):} $\langle \psi | (c_a c_b)^\dagger c_c c_d | \psi \rangle \ge 0 \implies Q(\rho^{(2)}) \succeq 0$.
    \item \textbf{G-Condition (Particle-Hole Positivity):} $\langle \psi | (c_a^\dagger c_b)(c_c^\dagger c_d)^\dagger | \psi \rangle \ge 0 \implies G(\rho^{(2)}) \succeq 0$.
\end{enumerate}
These matrix inequalities define a convex spectrahedron $\mathcal{S}_{\mathrm{SDP}}^{(2)}$ that encloses the convex hull of all physical states. Minimizing $\Tr(H \rho^{(2)})$ over this SDP spectrahedron yields a rigorous lower bound:
\begin{equation}
\label{eq:m1_1b_energy_sandwich}
E_{\mathrm{SDP}}^{(2)} \le E_0 \le E_{\mathrm{MF}}.
\end{equation}
The two-site covariance positivity condition $\mathcal{C}_{ij} \succeq 0$ proven in Chapter~\ref{ch:m1_4_5} is the exact local two-site manifestation of this semidefinite hierarchy.

% EDITORIAL QUESTION: How do quantum-information metrics quantify state distance and non-Gaussianity?
\section{Quantum Information Diagnostics: Entropy, Metrics, and Non-Gaussianity}
\label{sec:m1_1b_quantum_information}
% CLAIM: CLM-M1-QF-007 (provenance: doc/rnd/phase_1/m1_5_t3_quasifree_states_wick_variational.md, test: test/unit/test_quasifree_states.py::test_tst_qf_05_gaussian_entropy_and_maximum_entropy)

\begin{definition}[Gaussian von Neumann Entropy]
\label{def:m1_1b_gaussian_entropy}
Let $\Omega \in \mathfrak{so}(2n)$ have symplectic eigenvalues $\nu_k \in [0, 1]$ ($k = 1, \dots, n$). The von Neumann entropy of the corresponding thermal quasi-free state is:
\begin{equation}
\label{eq:m1_1b_von_neumann_entropy}
S(\omega_\Omega) = \sum_{k=1}^n h\left(\frac{1 + \nu_k}{2}\right), \qquad h(p) \equiv -p \log_2 p - (1-p) \log_2(1-p).
\end{equation}
By Jaynes' maximum entropy principle, the Gaussian state maximizes von Neumann entropy among all density matrices sharing the covariance matrix $\Omega$.
\end{definition}

\begin{definition}[Trace Distance and Uhlmann Fidelity]
\label{def:m1_1b_fidelity}
For two pure quasi-free states $\Omega_1, \Omega_2 \in \Gr(n, 2n)$, the Uhlmann fidelity evaluates to the closed-form determinant:
\begin{equation}
\label{eq:m1_1b_uhlmann_fidelity}
F(\Omega_1, \Omega_2) = |\langle \Psi_1 | \Psi_2 \rangle|^2 = \det\left(\frac{\Id - \Omega_1 \Omega_2}{2}\right)^{1/2}.
\end{equation}
\end{definition}

\begin{definition}[Non-Gaussianity Quantum Relative Entropy]
\label{def:m1_1b_non_gaussianity}
For any true many-body ground state $\rho_{\mathrm{exact}}$, its non-Gaussianity is measured by the quantum relative entropy to its Gaussian surrogate $\omega_{\mathrm{G}}[\rho]$ having the identical two-point covariance:
\begin{equation}
\label{eq:m1_1b_relative_entropy}
D_{\mathrm{NG}}(\rho) \equiv S(\rho \| \omega_{\mathrm{G}}[\rho]) = S(\omega_{\mathrm{G}}[\rho]) - S(\rho) \ge 0.
\end{equation}
By Pinsker's inequality, $D_{\mathrm{NG}}(\rho)$ bounds the trace distance to the Gaussian manifold: $\|\rho - \omega_{\mathrm{G}}\|_1 \le \sqrt{2 \ln 2 \cdot D_{\mathrm{NG}}(\rho)}$.
\end{definition}

% EDITORIAL QUESTION: How does this model-independent theory specialize to translation-invariant honeycomb lattices?
\section{Synthesis and Bridge to the Model-Specific Bloch Representation}
\label{sec:m1_1b_bridge_to_model}
The general theory developed in this chapter establishes that:
\begin{enumerate}
    \item Wick's theorem is a rigorous property of quasi-free states on $\mathcal{A}_{\mathrm{CAR}}$, not an operator approximation.
    \item The mean-field Hamiltonian $h(\Omega) = 4 \frac{\partial \mathcal{E}}{\partial \Omega}$ is the functional gradient on the real orthogonal Grassmannian $\Gr(n, 2n)$.
    \item Scalar double-counting subtraction $C[\Omega] = E_{\mathrm{int}}(\Omega)$ is the exact consequence of Euler's homogeneous function theorem on quadratic forms.
    \item Local gauge invariance is enforced by KKT stationarity $\omega_\Omega(\tau_i^\gamma) = 0$, leading to the variational upper bound and the Gutzwiller projection gap.
\end{enumerate}

In the subsequent chapter (\Cref{ch:m1_2}), we apply this universal framework to the \textbf{extended Kitaev honeycomb model}. Exploiting discrete translation symmetry, the global Grassmannian decomposes into a continuous direct product of momentum-space fibers:
\begin{equation}
\label{eq:m1_1b_grassmannian_fiber_bundle}
\Gr(4N, 8N) \xrightarrow{\text{Fourier}} \prod_{\mathbf{k} \in \mathrm{BZ}} \Gr(4, 8)_{\mathbf{k}},
\end{equation}
mapping the global matrix $h(\Omega)$ to the $8 \times 8$ Bloch Hamiltonian $\mathcal{H}_2(\mathbf{k})$.
